\section{Introduction}\label{sec:introduction}

The index of a log Calabi--Yau pair measures the degree in which its
rationally trivial log canonical class becomes an actual trivial line
bundle. A uniform index theorem asks whether that degree can be chosen from
the dimension and the boundary coefficients alone. For a nonnormal pair,
this question includes a descent problem: trivializations on the components
of the normalization need not agree along the double locus.

We work over an algebraically closed field $k$ of characteristic zero.
A scheme is \emph{demi-normal} if it is reduced, satisfies Serre's condition
$S_2$, is seminormal, and has only nodes in codimension one. Let $X$ be a
projective equidimensional demi-normal scheme, and let $B\ge0$ be a rational
Weil divisor with no component in the nonnormal locus. Write
$\nu:\bar X\to X$ for the normalization and $\bar C$ for its conductor
divisor. The pair $(X,B)$ is \emph{semi-log-canonical}, abbreviated slc,
if $K_X+B$ is rational Cartier and
\[
                 (\bar X,\bar C+\nu_*^{-1}B)
\]
is log canonical. These are the singularities that allow normal components
to meet along a double locus while retaining adjunction and pluricanonical
descent \citep{Kollar2013}.
We call $(X,B)$ a log Calabi--Yau pair when, in addition,
$K_X+B\sim_{\Q}0$: some positive Cartier multiple has trivial associated
line bundle. The \emph{index} of the pair is the least such positive
integer.

\begin{theorem}\label{thm:main}
For every integer $d\ge4$ and every finite set
$\Phi\subset[0,1]\cap\Q$, there is an integer $a(d,\Phi)>0$ with the
following property. Let $(X,B)$ be a connected projective equidimensional
semi-log-canonical log Calabi--Yau pair of dimension $d$ over $k$, with
every nonzero coefficient of $B$ in $\Phi$. Then
\[
            a(d,\Phi)(K_X+B)\text{ is Cartier},
       \qquad \OO_X\bigl(a(d,\Phi)(K_X+B)\bigr)\simeq\OO_X.
\]
\end{theorem}

The bound is independent of the number of irreducible components of $X$
and of the number of log canonical strata on their birational models.
It is a common trivializing multiple, so it clears the local Cartier
indices and kills the remaining global torsion at the same time.
Together with the theorem in dimensions at most three of
\citet[Corollary~1.6]{JiangLiu}, this proves the finite-rational-coefficient
slc index conjecture \citep[Conjecture~1.5]{JiangLiu}.

\subsection{History and relation to earlier work}
The distinction between existence of an index and a uniform index is
already visible in numerical-dimension-zero abundance. Gongyo proved that
a numerically trivial log canonical divisor on a projective slc pair is
rationally linearly trivial \citep[Theorem~1.5]{Gongyo2013}; the resulting
trivializing multiple may depend on the pair. The index conjecture asks
for a common multiple after fixing the dimension and the boundary
coefficients. Its finite-rational-coefficient slc formulation is stated,
for example, in \citet[Conjecture~1.5]{JiangLiu}. The nonnormal setting is
important because slc pairs occur as limits in moduli problems; see
\citet{Kollar2013} and the boundedness theory for polarized slc
Calabi--Yau pairs in \citet{Birkar2023}.

Pluricanonical representations have long connected index problems to
descent. Fujino developed admissible sections in his abundance theorem
for slc threefolds \citep{Fujino2000}, and formulated uniform boundedness
of canonical representations for canonical varieties with trivial canonical
divisor \citep[Conjecture~3.2]{Fujino2001}. His local-index results and
Ishii's global-index theorem for non-klt log canonical threefold pairs
with standard coefficients made this connection explicit
\citep{Fujino2001,Ishii2000}. Here the standard coefficients are
$1$ and $1-1/q$ for positive integers $q$.
Fujino--Gongyo subsequently proved finiteness of the log pluricanonical
representation for every projective lc pair with semiample log canonical
divisor \citep[Theorem~1.1]{FujinoGongyo2014}. Uniform boundedness of these
finite images is the additional issue in the index problem.

Xu organized reductions among the normal index conjecture, its slc
version, and boundedness of pluricanonical representations
\citep{Xu2019}. In particular, his Theorem~1.11 gives the implication
from normal indices in dimension $d$ and bounded representations in
dimension $d-1$ to slc indices in dimension $d$.
Jiang proved the global index bound for boundary-free klt threefolds
\citep[Corollary~1.7]{Jiang2021}.
Jiang--Liu bounded the representations of lc surfaces with a fixed
trivializing multiple, obtaining the slc index theorem in dimensions at
most three and the non-klt lc case in dimension four
\citep[Theorem~1.4 and Corollaries~1.6--1.7]{JiangLiu}.
Xu also treated klt fourfold pairs with nonzero boundary
\citep[Theorem~1.14]{Xu2019}.

Several more recent results give bounds under different geometric
hypotheses. Birkar proved the index conjecture for rationally connected
klt pairs in every dimension, allowing a fixed rational DCC coefficient
set \citep[Corollary~1.8]{Birkar2025}.
Filipazzi--Mauri--Moraga obtained sharp bounds for coregularity-zero
pairs, including slc pairs \citep{FMM2025}; here coregularity zero means
that a dlt model of the normalization with its induced boundary has a
zero-dimensional minimal lc stratum.
Figueroa--Filipazzi--Moraga--Peng developed corresponding index
and complement results in low coregularity \citep{FFMP2025}.
Masamura proved boundedness of the indices of smooth Calabi--Yau
fourfolds \citep[Corollary~1.6]{Masamura2025}.
Uniform bounds need not be numerically small: Singh constructed smooth
Calabi--Yau varieties whose indices grow doubly exponentially with
dimension \citep{Singh2026}.

The proof here follows the established normal-to-slc strategy. It uses
the uniform normal log canonical index theorem of the companion
\emph{Uniform Pluricanonical Iitaka Fibrations}
\citep[Theorem~1.2]{UPI}, together with the all-dimensional good-model
theorem of \emph{Log abundance in characteristic zero}
\citep[Theorem~11.1]{LA}. These are substantial inputs.
Integral cohomology already bounds pluricanonical characters in terms of
middle Betti numbers of suitable covers \citep[Lemma~2.7]{JiangLiu}.
The additional uniformity comes from a rank bound for the rational Hodge
structure generated by a holomorphic top form and the resulting character
bound for arbitrary integral klt pairs with fixed index. The latter is
the klt case of the bounded-representation problem formulated in
\citet[Conjecture~1.1]{JiangLiu}. The final comparison of conductor
residues uses Koll{\'a}r's linking theorem and its even-degree residue
compatibility \citep[Theorem~10 and Proposition~14]{KollarSources}.

The rank argument combines local positivity, moving subvarieties, and
section counts. Its order propagation belongs to the
Ein--K{\"u}chle--Lazarsfeld method \citep{EKL1995}; we use the precise
chain and scalar constructions of \citet[Sections~5--6]{U4} and
\citet[Section~6]{UPI}. Cotangent semipositivity
\citep{GKP2016,CP2019}, divisorial Zariski decomposition
\citep{Nakayama2004,Boucksom2004,BFJ2009}, and effective birationality
\citep{Birkar2023} supply the positivity estimates.
The character argument builds on Matsumura--Wang's decomposition of klt
log Calabi--Yau pairs \citep{MatsumuraWang2025} and the singular
Beauville--Bogomolov theory
\citep{GKP2016,Druel2018,GGK2019,HoringPeternell2019}.
For rationally connected factors we use Han--Jiang's boundedness theorem
and Jiang--Liu's bounded-family representation theorem
\citep{HanJiang,JiangLiu}. The proofs below isolate the new uniform
estimates and explain the adaptations of these earlier constructions.


\subsection{The character obstruction and the proof}

We first describe the descent problem that dictates the proof.
The uniform normal index theorem of \citet[Theorem~1.2]{UPI} supplies
one even integer $m$, depending only on $d$ and $\Phi$, for all the
normalization components. On the $i$th component its trivialization can
be written as a rational $m$-canonical form $\theta_i$ with the prescribed
logarithmic poles. On a crepant divisorial log terminal model, logarithmic
residue restricts this form to the conductor strata. At a generic node,
the two residues differ by a nonzero scalar $r_e$; the existence of some
global trivializing multiple ensures that the ratio is constant.

Make a finite graph whose vertices are normalization components and whose
edges are the generic nodes. Rescaling $\theta_i$ changes its incident
edge ratios. Thus simultaneous matching amounts to making every product
of ratios around a closed walk equal to one. A uniform power of all the
forms will accomplish this if the orders of these products are uniformly
bounded.

Minimal log canonical strata supply the necessary interpretation of a
closed walk. Their adjoint pairs are klt. Inside one normalization
component, $\PP^1$-linking identifies any two minimal strata by a crepant
birational map that preserves the even-degree residue. Across an edge,
we prove a corresponding birational residue comparison with multiplier
$r_e$. A closed walk therefore induces a crepant birational self-map of
one minimal klt stratum, whose action on its trivializing form is exactly
the product of the edge ratios. Section~\ref{sec:residues} carries out
these comparisons, including nonsplit nodes and the final $S_2$ extension.

The main uniformity problem is consequently the following one. For a
projective integral klt pair $(V,\Delta)$ satisfying
$m(K_V+\Delta)\sim0$, bound, in terms of $\dim V$ and $m$, the order of
the scalar by which a crepant birational self-map acts on a trivializing
$m$-canonical form. We prove this in Theorem~\ref{thm:characters}.
The structural decomposition of such pairs separates a rationally
connected factor, an abelian factor, and irreducible Calabi--Yau or
symplectic factors. The rationally connected factor is controlled by
boundedness and log pluricanonical representations. For the other factors,
the character appears in integral middle cohomology.

The relevant cohomology is the smallest rational Hodge substructure
$T(U)\subset H^n(U,\Q)$ containing the holomorphic top forms of a smooth
projective resolution $U$. It is birationally invariant and carries a
polarized integral lattice. A bound for its rank therefore bounds the
degree of the cyclotomic polynomial of a top-form character. Theorem~\ref{thm:rank}
provides that rank bound under the rational-image hypotheses satisfied
by the irreducible factors.

Two parts of the rank argument are useful independently. First, a
dimension-dependent estimate for the variation of a Hodge norm gives a
strict improvement over the largest possible metric growth whenever a
specified top cup product vanishes. Applied to a class on the blowup of
the diagonal, it bounds $\dim T(U)$ from a normalized volume and a lower
Seshadri bound. These arguments occupy Sections~\ref{sec:hodge}
and~\ref{sec:diagonal}. Second, Section~\ref{sec:rank} removes the
Seshadri hypothesis. Small-order moving subvarieties determine rational
fibrations; a maximal such fibration reduces the issue to vertical jets
and horizontal section counts. A calibrated big divisor keeps the two
counts on the same scale, and a finite power-map argument makes the
horizontal estimate linear in the sheaf rank.

Section~\ref{sec:characters} converts this rank bound into the required
uniform character exponent. Section~\ref{sec:residues} then proves
Theorem~\ref{thm:main}. The substantial companion inputs---the normal
index theorem, good minimal models, and the scalar and structural
results used in these reductions---are stated with their precise roles
in Section~\ref{sec:preliminaries} and at their points of use.
